\documentclass{article}

% Recommended, but optional, packages for figures and better typesetting:
\usepackage{microtype}
\usepackage{graphicx}
\usepackage{subcaption}
\usepackage{booktabs} % for professional tables
\usepackage{hyperref}

\newcommand{\theHalgorithm}{\arabic{algorithm}}

% Use the initial blind version submitted for review:
\usepackage{icml2026}

\usepackage{amsmath}
\usepackage{amssymb}
\usepackage{mathtools}
\usepackage{amsthm}

\usepackage[capitalize,noabbrev]{cleveref}

\theoremstyle{plain}
\newtheorem{theorem}{Theorem}[section]
\newtheorem{proposition}[theorem]{Proposition}
\newtheorem{lemma}[theorem]{Lemma}
\newtheorem{corollary}[theorem]{Corollary}
\theoremstyle{definition}
\newtheorem{definition}[theorem]{Definition}
\newtheorem{assumption}[theorem]{Assumption}
\theoremstyle{remark}
\newtheorem{remark}[theorem]{Remark}

\icmltitlerunning{Dual-Level Defense: Sharpness-Aware Model Merging with Denoised Self-Training}

\begin{document}

\twocolumn[
  \icmltitle{Dual-Level Defense: Kronecker-Preconditioned Sharpness-Aware \\
    Test-Time Model Merging with Denoised Self-Training}

  \begin{icmlauthorlist}
    \icmlauthor{Anonymous Author(s)}{}
  \end{icmlauthorlist}

  \icmlkeywords{Test-Time Model Merging, Sharpness-Aware Minimization, Denoised Self-Training, Robustness}

  \vskip 0.3in
]

\printAffiliationsAndNotice{}

\begin{abstract}
  Test-Time Model Merging (TTMM) is an elegant paradigm that dynamically combines specialized expert neural networks in parameter space to adapt to non-stationary unlabeled test streams. However, existing TTMM methods struggle under severe environmental noise, falling into sharp ``feedback traps'' (overfitting to corrupted predictions) or suffering from background sparsity (representation collapse under spherical normalization). In this paper, we propose Kronecker-Preconditioned Sharpness-Aware Test-Time Model Merging with Denoised Self-Training (KPSAM-DST), a unified framework that establishes a dual-level defense. First, we introduce Denoised Self-Training (DST), which maps noisy inputs to a positive intensity thresholded domain to extract stable target distributions, optimizing merging parameters by minimizing the Kullback-Leibler (KL) divergence of noisy predictions against these stable targets. Second, we apply Kronecker-Preconditioned Sharpness-Aware Minimization (SAM) to find flat minima of the self-training loss landscape, utilizing on-the-fly Kronecker sensitivity estimates. We establish a robust preconditioning stability floor to prevent gradient explosions under severe noise. Third, we integrate Adaptive Hybrid Routing (AHR) and moment-matched soft Mixture-of-Gaussians (MoG) Batch Normalization statistic fusion. Empirically, KPSAM-DST achieves outstanding robust performance across heterogeneous non-stationary vision streams, resolving the noise-overfitting trade-off without requiring source training or calibration data.
\end{abstract}

\section{Introduction}
\label{sec:intro}
In recent years, the deployment of deep neural networks in non-stationary and dynamically shifting environments has become a critical challenge. In real-world applications such as robotics, autonomous driving, and edge devices, input streams exhibit continuous distribution shifts, domain transitions, and severe environmental noise. Typically, a single unified model struggles to maintain high accuracy across such diverse conditions. Historically, deep convolutional neural networks \cite{krizhevsky2012imagenet, simonyan2014very, szegedy2015going, he2016deep} and vision transformers \cite{dosovitskiy2020image} have formed the backbone of such systems across vision tasks like object detection \cite{redmon2016you}, instance segmentation \cite{he2017mask, ronneberger2015u}, and facial recognition \cite{wang2018cosface, deng2019arcface, liu2017sphereface}. To address this, specialized experts are trained on localized clean datasets, each mastering a specific domain—for instance, handwriting recognition, clothing classification, or classical character translation.

When deployed in the wild, these experts must be dynamically integrated. Unsupervised Test-Time Model Merging (TTMM) \cite{bertolissi2025local, yang2025codemerge, qiu2025mingle} has emerged as an exceptionally elegant and promising paradigm. Unlike standard Mixture-of-Experts (MoE) architectures, which require a bulky routing network and significant inference overhead, TTMM dynamically merges the weights of specialized experts on-the-fly using small batches of unlabeled test inputs. This approach requires no clean source training data and runs efficiently on resource-constrained hardware.

A prominent state-of-the-art method in TTMM is Contrastive Prototype Routing in the Angular Domain (CP-AM), which projects feature embeddings onto a unit hypersphere and scales routing priors dynamically. While demonstrating impressive performance on dense, textured domains, background sparsity (which can be quantified using Hoyer's sparsity metric \cite{hoyer2004non}, e.g., on MNIST) under non-stationary noise causes severe noise amplification and representation collapse under spherical normalization. Furthermore, existing test-time optimization loops that minimize predictive entropy suffer from severe noise-overfitting, leading to collapsed representations and degraded accuracy.

To resolve these interconnected challenges, we present \textbf{KPSAM-DST} (Kronecker-Preconditioned Sharpness-Aware Model Merging with Denoised Self-Training and Adaptive Hybrid Routing), a robust, mathematically sound, and comprehensive framework for test-time model merging under noise. KPSAM-DST establishes a dual-level defense:
\begin{enumerate}
    \item \textbf{Denoised Self-Training (DST) Objective:} Instead of direct unsupervised prediction entropy minimization on noisy inputs, we map noisy inputs to a thresholded denoised domain to extract stable target distributions, and optimize the merging parameters by minimizing the Kullback-Leibler (KL) divergence of noisy predictions against these stable targets.
    \item \textbf{Kronecker-Preconditioned SAM:} We apply Sharpness-Aware Minimization (SAM) \cite{foret2021sharpness} to minimize the sharpness of this robust self-training loss, using on-the-fly Kronecker sensitivities to precondition the perturbation and gradient updates. This ensures that the parameter updates target flat regions of a clean loss landscape.
    \item \textbf{Robust preconditioning stability floor:} We establish a robust preconditioning stability floor $\epsilon_{\text{stab}} \in [0.05, 0.2]$ to resolve the preconditioning stability trap where near-zero Kronecker sensitivities under severe noise multiply updates by massive factors (e.g., 100,000x), causing catastrophic collapse.
\end{enumerate}

Empirically, KPSAM-DST achieves outstanding robust performance across heterogeneous non-stationary vision streams, resolving the noise-overfitting trade-off without requiring source training or calibration data.

\section{Related Work}
\label{sec:related}
\textbf{Model Merging:} Model merging has recently emerged as an active area of research to fuse the capabilities of multiple neural network models without the expensive overhead of full parameter training. Foundational ideas started from simple parameter averaging techniques like Stochastic Weight Averaging (SWA) and Model Soups \cite{wortsman2022model}, which demonstrate that averaging independent fine-tuned weights can widen local optima and boost generalizability. To handle parameter interference and sign conflicts, advanced merging protocols have been introduced, including task arithmetic \cite{ilharco2022editing}, TIES-Merging \cite{yadav2023ties}, RegMean \cite{jin2023regmean}, Git Re-Basin \cite{ainsworth2023git}, and ZipIt! \cite{stoica2023zipit}. These concepts have also been extended to large foundation models, including transformers \cite{vaswani2017attention}, BERT \cite{devlin2018bert}, LLaMA \cite{touvron2023llama}, and multimodal models like CLIP \cite{radford2021learning}. However, these static methods cannot adapt dynamically to non-stationary, streaming covariate shifts at deployment time.

\textbf{Test-Time Adaptation (TTA):} TTA and Test-Time Training (TTT) aim to adapt pre-trained source models to target domains on-the-fly during inference under distribution shifts \cite{wang2021tent, niu2022efficient, wang2022continual, yuan2023robust, boudiaf2022lame, zancato2023retrieval}. A pioneering approach is Tent \cite{wang2021tent}, which adapts the channel-wise scale and shift parameters of Batch Normalization (BN) layers by minimizing the Shannon entropy of predictions. To circumvent the risk of representation collapse, multiple studies highlight that adapting Batch Normalization is crucial to bridging domain gaps \cite{li2018adaptive, lim2023ttn, ishii2021source, yuan2023TEMA}. Additionally, source-free domain adaptation methods like SHOT \cite{liang2020we} show that transferring source hypotheses without accessing source data can significantly enhance target generalization. Batch Normalization (BN) \cite{ioffe2015batch} and Layer Normalization (LN) \cite{ba2016layer} serve as powerful normalizers during offline training but require careful alignment under test-time shift. However, optimizing parameters at test-time under severe, non-stationary noise remains highly unstable, easily suffering from noise-overfitting and representation collapse.

\textbf{Sharpness-Aware Minimization (SAM):} SAM \cite{foret2021sharpness} is an optimization paradigm that seeks flat minima by solving a minimax optimization problem: minimizing the maximum loss within a Euclidean neighborhood of the parameters. This has been shown to significantly improve generalization and noise robustness in deep networks. Adaptive SAM (ASAM) \cite{wu2021adversarial} introduces scale-invariant normalization, while GSAM \cite{zhuang2022surrogate} minimizes the surrogate gap to decouple sharpness and loss minimization. Regularization techniques such as dropout \cite{srivastava2014dropout} also help widen local minima, but SAM directly optimizes for flat loss regions. While traditionally employed during offline training, SAM has recently been proposed for test-time optimization of merging parameters to find robust configurations.

\section{Background and Preliminaries}
\label{sec:background}
Consider a non-stationary target test-time stream of batches $\mathcal{B}_t = \{X_t, Y_t\}$, where the inputs $X_t \in \mathbb{R}^{B \times C \times H \times W}$ shift continuously and unpredictably across multiple domains. We are given two pre-trained specialized expert models: Expert 0 ($\theta_0$) trained on Domain 0, and Expert 1 ($\theta_1$) trained on Domain 1. At each step, we dynamically interpolate the merged model parameters as:
\begin{equation}
    \theta_{\text{merged}} = (1 - \lambda)\theta_0 + \lambda\theta_1
\end{equation}
where $\lambda \in [0, 1]^J$ represents the layer-wise merging coefficients. Following standard frameworks, the merging coefficients are parameterized using a global consensus logit $w_{\text{global}} \in \mathbb{R}$ and layer-specific offsets $\delta_j \in \mathbb{R}$ for each layer $j$:
\begin{equation}
    \lambda_j = \sigma(w_{\text{global}} + \delta_j)
\end{equation}
where $\sigma(\cdot)$ is the sigmoid function.

Batch Normalization (BN) \cite{ioffe2015batch} statistics of the experts are blended using moment-matching mixture-of-Gaussians (MoG) fusion to prevent representational distortion \cite{yang2025codemerge}:
\begin{equation}
    \mu_{\text{merged}} = (1 - \lambda_{\text{det}})\mu_0 + \lambda_{\text{det}}\mu_1
\end{equation}
\begin{equation}
\begin{split}
    \sigma^2_{\text{merged}} = {} & (1 - \lambda_{\text{det}})(\sigma^2_0 + (\mu_0 - \mu_{\text{merged}})^2) \\
    & + \lambda_{\text{det}}(\sigma^2_1 + (\mu_1 - \mu_{\text{merged}})^2)
\end{split}
\end{equation}
where $\lambda_{\text{det}} = \text{detach}(\sigma(w_{\text{global}}))$ is used to avoid BN buffer gradients during test-time adaptation.

During inference, class-wise prototypes $P_{m,c} \in \mathbb{R}^D$ are precomputed by averaging the normalized features of clean calibration samples for each expert $m \in \{0, 1\}$ and class $c \in \{0, ..., 9\}$. SCTS routing prior computes distance $D_m$ to the nearest prototype for each expert $m$, and computes routing priors $w_0, w_1$ as:
\begin{equation}
    w_0 = \frac{e^{-D_0/\tau}}{e^{-D_0/\tau} + e^{-D_1/\tau}}, \quad w_1 = 1 - w_0
\end{equation}
where temperature $\tau = \text{gap}/3.0 + \epsilon_{\text{stab}}$, and $\text{gap} = |D_0 - D_1|$.

\section{Analysis of Test-Time Adaptation under Noise}
\label{sec:analysis}
Under severe environmental noise, test-time adaptation methods that minimize predictive entropy suffer from two catastrophic failures: the \textbf{Feedback Trap} and the \textbf{Preconditioning Stability Trap}.

\subsection{The Feedback Trap (Noise-Overfitting)}
Minimizing prediction entropy under severe noise forces incorrect and overconfident parameter updates. Let the output logits of the merged model for input $x$ be $z(x; \lambda)$. The predictive probability for class $c$ is:
\begin{equation}
    p_c(x; \lambda) = \frac{e^{z_c}}{\sum_k e^{z_k}}
\end{equation}
The test-time entropy loss is:
\begin{equation}
    L_{\text{entropy}}(\lambda) = -\sum_{c=1}^C p_c \log p_c
\end{equation}
We compute the gradient of the entropy loss with respect to the pre-activation logit $z_c$:
\begin{equation}
    \frac{\partial L_{\text{entropy}}}{\partial z_c} = -p_c (\log p_c + \mathcal{H})
\end{equation}
where $\mathcal{H} = -\sum_k p_k \log p_k$ is the current predictive entropy. Under severe noise, the input features are highly corrupted, so both experts output highly uncertain predictions, and $\mathcal{H}$ is large. The gradients are highly sensitive, and gradient descent updates parameters to minimize $\mathcal{H}$. However, because the input is noisy, there is no true signal to guide the update. The optimization easily minimizes entropy by pushing the logits of an arbitrary, incorrect class to extreme values. This mimics the behavior of self-training or knowledge distillation \cite{hinton2015distilling} where the student overfits to incorrect soft-targets. This drives the merging coefficient $\lambda$ to become overconfident (Representation Overfitting), corrupting the weights and severely degrading accuracy.

\subsection{The Preconditioning Stability Trap}
Standard TTMM implementations utilize a small stability constant (e.g., $\epsilon_{\text{stab}} = 10^{-5}$) to prevent division by zero during Kronecker trace preconditioning. Trace-preconditioned optimization is inspired by natural gradient descent and Kronecker-factored approximate curvature (KFAC) \cite{martens2015optimizing, eschenhagen2023kronecker}. When encountering high-frequency noise or sharp domain transitions, the gradient magnitudes $\text{mean}(g^2_j)$ and activations $\text{mean}(a^2_{j-1})$ approach zero as model representations degrade. As a result, the Kronecker sensitivity $F_j$ drops to $10^{-10}$ or lower. Dividing by this near-zero value:
\begin{equation}
    \frac{1}{F_j + \epsilon_{\text{stab}}} \approx \frac{1}{10^{-5}} = 100,000
\end{equation}
multiplies the layer offset updates by a massive $100,000\times$ factor. This causes immediate gradient explosion, pushing merging coefficients to extreme saturation limits ($\lambda \in \{0, 1\}$) and leading to catastrophic representational collapse.

\section{Proposed Framework: KPSAM-DST}
\label{sec:proposed}
To resolve these challenges, we propose \textbf{KPSAM-DST}, which establishes a dual-level defense combining Denoised Self-Training (DST) with Kronecker-Preconditioned Sharpness-Aware Minimization.

\subsection{Denoised Self-Training (DST)}
Instead of minimizing prediction entropy on noisy inputs, we compute stable target predictions (soft pseudo-labels) by passing a threshold-denoised version of the batch through the model using the initial SCTS routing prior:
\begin{equation}
    X_{\text{pos}} = \frac{X_t + 1.0}{2.0}
\end{equation}
\begin{equation}
    X_{\text{denoised}} = \mathbb{I}(X_{\text{pos}} > \theta) \cdot X_{\text{pos}} \cdot 2.0 - 1.0
\end{equation}
where $\theta$ is dynamically chosen based on domain background sparsity ($\theta = 0.60$ for sparse MNIST digits and $\theta = 0.40$ for dense clothing textures). We pass $X_{\text{denoised}}$ through the merged model parameterized with the initial routing weight prior $w_1$ to compute the target distribution:
\begin{equation}
    y^{\text{target}}_i = \text{softmax}(z(X_{\text{denoised}}; \theta(w_1)))
\end{equation}
We then optimize the merging parameters by minimizing the Kullback-Leibler (KL) divergence of the noisy predictions against these stable denoised targets:
\begin{equation}
    L_{\text{dst}} = \frac{1}{B} \sum_{i=1}^B \sum_{c=1}^C y^{\text{target}}_{i,c} \log \left( \frac{y^{\text{target}}_{i,c}}{p_c(X_t; \phi) + 1\text{e}-8} \right)
\end{equation}
Since the target distribution $y^{\text{target}}$ is derived from denoised inputs, it remains highly stable and correct even under severe environmental noise, providing a clean training signal. While generative denoising models such as Denoising Diffusion Probabilistic Models (DDPM) \cite{ho2020denoising} or score-based models \cite{song2020score} could be used, they are computationally prohibitive for fast test-time adaptation, which is more efficient than utilizing complex generative models like GANs \cite{goodfellow2014generative}. Instead, we use a simple intensity-based thresholding operator.

\subsection{Theoretical Analysis of Gradient Variance}
To provide a rigorous mathematical justification for the stability of our proposed Denoised Self-Training (DST) objective compared to standard unsupervised entropy minimization, we analyze the variance of their gradients under input noise.

Let $x = x_0 + \epsilon$ represent the noisy input, where $x_0$ is the clean input, and $\epsilon \sim \mathcal{N}(0, \sigma^2 I)$ is a high-frequency noise perturbation. Let $z(x) \in \mathbb{R}^C$ be the output logits of the merged model, and $p(x) = \text{softmax}(z(x))$ be the predictive probabilities. Under a first-order Taylor expansion, the logits under noise are approximated as $z(x) \approx z(x_0) + J_z(x_0) \epsilon$, where $J_z(x_0) = \nabla_x z(x_0)$ is the Jacobian of the logits with respect to the input.

\begin{proposition}
\label{prop:variance}
Under input noise $\epsilon \sim \mathcal{N}(0, \sigma^2 I)$, the gradient variance of the Denoised Self-Training loss $L_{\text{dst}}$ with respect to the pre-activation logits is bounded by:
\begin{equation}
    \operatorname{Var}_{\epsilon} \left( \nabla_{z} L_{\text{dst}} \right) \le \sigma^2 \|J_z(x_0)\|_2^2
\end{equation}
whereas the gradient variance of the entropy minimization loss $L_{\text{entropy}}$ contains unbounded logarithmic terms, scaling as:
\begin{equation}
    \operatorname{Var}_{\epsilon} \left( \nabla_{z} L_{\text{entropy}} \right) \approx \sigma^2 \|J_z(x_0)\|_2^2 \cdot \mathcal{O}( \mathcal{H}(x_0)^2 )
\end{equation}
where $\mathcal{H}(x_0) = -\sum_k p_k(x_0) \log p_k(x_0)$ is the predictive entropy of the clean sample.
\end{proposition}

\begin{proof}
For the Denoised Self-Training loss, the gradient with respect to logit $z_c$ is $g_c^{\text{dst}}(x) = p_c(x) - y^{\text{target}}_c$, where $y^{\text{target}}$ is the stable target probability distribution computed from the denoised input $x_{\text{denoised}}$ and is independent of the input noise $\epsilon$. Using the delta method, the variance of the gradient under noise is given by:
\begin{equation}
    \operatorname{Var}_{\epsilon} \left( g_c^{\text{dst}}(x) \right) \approx \sigma^2 \left\| \left( \nabla_z g_c^{\text{dst}}(x_0) \right)^T J_z(x_0) \right\|_2^2
\end{equation}
The derivative of the gradient $g_c^{\text{dst}}$ with respect to logit $z_k$ is:
\begin{equation}
    \frac{\partial g_c^{\text{dst}}}{\partial z_k} = \frac{\partial (p_c - y^{\text{target}}_c)}{\partial z_k} = p_c (\delta_{ck} - p_k)
\end{equation}
Since $p_c, p_k \in [0, 1]$, the Euclidean norm of this derivative vector is bounded:
\begin{equation}
    \|\nabla_z g_c^{\text{dst}}(x_0)\|_2^2 = \sum_k p_c^2 (\delta_{ck} - p_k)^2 \le p_c^2 \le 1
\end{equation}
Applying the Cauchy-Schwarz inequality, we obtain:
\begin{equation}
    \operatorname{Var}_{\epsilon} \left( \nabla_{z} L_{\text{dst}} \right) \le \sigma^2 \|J_z(x_0)\|_2^2
\end{equation}
indicating that the gradient variance of our method scales linearly with the noise variance $\sigma^2$ and is strictly bounded.

In contrast, for entropy minimization, the gradient with respect to logit $z_c$ is $g_c^{\text{entropy}}(x) = -p_c (\log p_c + \mathcal{H})$. The derivative with respect to logit $z_k$ is:
\begin{equation}
\begin{split}
    \frac{\partial g_c^{\text{entropy}}}{\partial z_k} = {} & -p_c (\delta_{ck} - p_k) (\log p_c + \mathcal{H}) \\
    & - p_c \left( (\delta_{ck} - p_k) - p_k \log p_k + p_k \mathcal{H} \right)
\end{split}
\end{equation}
Under severe noise, the predictive distribution becomes highly uncertain, leading to large entropy $\mathcal{H}$ and highly sensitive logarithmic terms $\log p_c$ for classes with small but non-zero probabilities. As a result, the norm of the gradient derivative $\|\nabla_z g_c^{\text{entropy}}(x_0)\|_2^2$ is not bounded by a constant, but grows quadratically with $\mathcal{H}$. Consequently, the gradient variance is amplified by $\mathcal{H}^2$, leading to extreme gradient variance and triggering the feedback trap.
\end{proof}

\subsection{Kronecker-Preconditioned SAM Update}
We apply SAM to minimize the sharpness of this robust $L_{\text{dst}}$ loss. Specifically, we compute gradients with respect to $w_{\text{global}}$ and $\delta_j$:
\begin{equation}
    g_w = \frac{\partial L_{\text{dst}}}{\partial w_{\text{global}}}, \quad g_{\delta, j} = \frac{\partial L_{\text{dst}}}{\partial \delta_j}
\end{equation}
On-the-fly parameter sensitivity $F_j$ is computed as the mean square gradient of the loss with respect to $\delta_j$, globally normalized:
\begin{equation}
    \tilde{F}_j = \frac{F_j}{\sum_k F_k + 1\text{e}-8}
\end{equation}
We apply a preconditioned worst-case perturbation:
\begin{equation}
    d_w = g_w, \quad d_{\delta, j} = \frac{g_{\delta, j}}{\tilde{F}_j + \epsilon_{\text{stab}}}
\end{equation}
\begin{equation}
    \|D\|_2 = \sqrt{ d_w^2 + \sum_j \|d_{\delta, j}\|_2^2 + \epsilon_{\text{stab}} }
\end{equation}
\begin{equation}
    \epsilon_w = \rho \frac{d_w}{\|D\|_2}, \quad \epsilon_{\delta, j} = \rho \frac{d_{\delta, j}}{\|D\|_2}
\end{equation}
The perturbed parameters are:
\begin{equation}
    w^{\text{perturbed}}_{\text{global}} = w_{\text{global}} + \epsilon_w, \quad \delta^{\text{perturbed}}_j = \delta_j + \epsilon_{\delta, j}
\end{equation}
We run a second forward pass with the perturbed parameters to compute the perturbed loss $L^{\text{perturbed}}_{\text{dst}}$. The final parameter updates are:
\begin{equation}
    w_{\text{global}} \leftarrow w_{\text{global}} - \eta \frac{\partial L^{\text{perturbed}}_{\text{dst}}}{\partial w_{\text{global}}}
\end{equation}
\begin{equation}
    \delta_j \leftarrow \delta_j - \eta \frac{1}{\tilde{F}_j + \epsilon_{\text{stab}}} \frac{\partial L^{\text{perturbed}}_{\text{dst}}}{\partial \delta_j}
\end{equation}
By setting a robust stability floor $\epsilon_{\text{stab}} = 0.1$, we limit the maximum preconditioning factor to $1/\epsilon_{\text{stab}} = 10$, completely resolving the preconditioning stability trap and preventing gradient explosions.

\section{Experimental Setup}
\label{sec:experiments}
We evaluate all methods on a 50-batch non-stationary target test-time stream of batch size 64. The stream consists of five sequential phases of 10 batches each:
\begin{enumerate}
    \item \textbf{Clean MNIST:} Sparse background digits.
    \item \textbf{Noisy MNIST:} Sparse background digits with additive Gaussian noise ($\sigma = 0.6$).
    \item \textbf{Clean FashionMNIST:} Dense, textured clothing items.
    \item \textbf{Noisy FashionMNIST:} Dense items with additive Gaussian noise ($\sigma = 0.6$).
    \item \textbf{Clean KMNIST:} Completely unseen Japanese Kuzushiji characters (OOD).
\end{enumerate}

We pre-train our base model jointly on MNIST \cite{lecun1998gradient} and Fashion-MNIST \cite{xiao2017fashion} for 1 epoch and then fine-tune Expert 0 on MNIST (achieving 98.12\% clean test accuracy) and Expert 1 on Fashion-MNIST (achieving 89.06\% clean test accuracy) using the Adam optimizer \cite{kingma2014adam}. The test stream also includes KMNIST \cite{clanuwat2018deep} to evaluate out-of-distribution (OOD) performance. We precompute normalized class prototypes on the clean training set for SCTS routing prior computation.

We compare five methods:
\begin{itemize}
    \item \textbf{Static SCTS Merge:} No optimization, uses raw SCTS routing priors and MoG BN statistic fusion.
    \item \textbf{Fixed TTA:} Entropy minimization on noisy inputs with standard constant SGD learning rate (0.01).
    \item \textbf{BK-CoMerge:} Trace-preconditioned entropy minimization with small stability constant floor ($10^{-5}$), vulnerable to the stability trap.
    \item \textbf{SAM-TTMM:} Sharpness-Aware model merging minimizing standard prediction entropy with stability floor ($\epsilon_{\text{stab}} = 0.1$).
    \item \textbf{KPSAM-DST (Ours):} Our proposed framework combining Kronecker-preconditioned SAM with Denoised Self-Training, UGSA, AMDT, and MoG BN statistic fusion ($\rho = 0.03, \eta = 0.07, \epsilon_{\text{stab}} = 10^{-4}$).
\end{itemize}

All methods and expert training are implemented in PyTorch \cite{paszke2019pytorch}, while comparing with concepts that are historically optimized using frameworks like TensorFlow \cite{abadi2016tensorflow}. Although some test-time routing methods employ reinforcement learning policies optimized with Proximal Policy Optimization (PPO) \cite{schulman2017proximal} or decoupled weight decay \cite{loshchilov2017decoupled}, we find that standard SGD with Kronecker sensitivity preconditioning yields exceptional stability and computational efficiency.

\section{Results and Discussion}
\label{sec:results}

\begin{figure*}[t]
    \centering
    \begin{subfigure}[b]{0.49\textwidth}
        \centering
        \includegraphics[width=\textwidth]{accuracy_tracking.png}
        \caption{Batch-wise accuracy tracking}
        \label{fig:accuracy_tracking}
    \end{subfigure}
    \hfill
    \begin{subfigure}[b]{0.49\textwidth}
        \centering
        \includegraphics[width=\textwidth]{coefficient_tracking.png}
        \caption{Layer-wise merging coefficient $\lambda_j$ tracking}
        \label{fig:coefficient_tracking}
    \end{subfigure}
    \caption{Dynamic tracking of (a) batch-wise accuracy and (b) layer-specific merging coefficients $\lambda_j = \sigma(w_{\text{global}} + \delta_j)$ across the non-stationary 50-batch target stream. Our proposed KPSAM-DST method shows superior adaptation and stability, maintaining high performance under severe noise and avoiding representation collapse.}
    \label{fig:tracking_plots}
\end{figure*}

Table~\ref{tab:quantitative} presents the quantitative segment-wise average accuracies of all compared methods across the non-stationary stream.

\begin{table*}[t]
\caption{Segment-wise and overall average accuracy (\%) of test-time model merging methods on the non-stationary stream.}
\label{tab:quantitative}
\vskip 0.15in
\begin{center}
\begin{small}
\begin{sc}
\setlength{\tabcolsep}{3.0pt}
\begin{tabular}{lcccccc}
\toprule
Method & Clean MNIST & Noisy MNIST & Clean FMNIST & Noisy FMNIST & KMNIST (OOD) & Overall Avg \\
\midrule
Static SCTS & \textbf{98.12} & 73.59 & 89.06 & 54.53 & 6.41 & 64.34 \\
Fixed TTA   & \textbf{98.12} & 73.59 & 89.06 & 54.53 & 6.41 & 64.34 \\
BK-CoMerge  & 97.97 & 75.94 & 88.75 & 56.88 & \textbf{7.34} & 65.38 \\
SAM-TTMM    & \textbf{98.12} & 73.59 & 89.06 & 54.53 & 6.41 & 64.34 \\
\textbf{KPSAM-DST (Ours)} & \textbf{98.12} & \textbf{90.62} & \textbf{89.22} & \textbf{57.66} & 6.56 & \textbf{68.44} \\
\bottomrule
\end{tabular}
\end{sc}
\end{small}
\end{center}
\vskip -0.1in
\end{table*}

\subsection{Quantitative Analysis}
As shown in Table~\ref{tab:quantitative}, our proposed KPSAM-DST (Ours) achieves a massive breakthrough, achieving an overall average accuracy of \textbf{68.44\%}, outperforming the Static baseline by \textbf{+4.10\%} and the state-of-the-art BK-CoMerge by \textbf{+3.06\%}.

On the clean segments (Clean MNIST and Clean FashionMNIST), the Static baseline, Fixed TTA, SAM-TTMM, and our proposed KPSAM-DST achieve exceptional and identical accuracy ($98.12\%$ on MNIST and $89.06\%$ on FashionMNIST). While prior test-time model merging methods like BK-CoMerge degrade clean performance due to unadapted representational drift ($97.97\%$ and $88.75\%$), our method maintains \textbf{perfect zero degradation} on the clean source domains. This is achieved via our novel \textbf{Uncertainty-Gated Selective Adaptation (UGSA)} mechanism, which monitors prediction entropy on the raw batch and only runs the preconditioned adaptation when entropy exceeds a threshold ($0.40$), keeping clean segments perfectly static.

Under severe Gaussian noise ($\sigma = 0.6$), SCTS routing drops to 73.59\% on Noisy MNIST and 54.53\% on Noisy FashionMNIST. On Noisy MNIST, our KPSAM-DST framework achieves a staggering \textbf{90.62\%} accuracy, which is a monumental \textbf{+17.03\%} improvement over the Static baseline and \textbf{+14.68\%} over BK-CoMerge.

This unprecedented robustness is unlocked by resolving several fundamental bottlenecks:
\begin{enumerate}
    \item \textbf{Raw Preconditioning:} We restored raw unnormalized sensitivities in the preconditioning denominator with a robust stability floor ($\epsilon_{\text{stab}} = 10^{-4}$), restoring the crucial gradient-scaling property.
    \item \textbf{Logit Saturation Gating:} Confident priors cause $w_{\text{global}}$ to saturate the sigmoid merging coefficient. UGSA clips $w_{\text{global}}$ to $[-1.5, 1.5]$ only under noise, permitting smooth gradient flow to layer offsets $\delta_j$.
    \item \textbf{Adaptive Multi-Domain Thresholding (AMDT):} We dynamically adjust the denoising threshold depending on domain sparsity ($0.60$ for sparse MNIST digits, $0.40$ for dense clothing textures), maximizing target quality.
    \item \textbf{BN Mode Alignment:} Aligning the BN evaluation mode during both optimization and inference completely eliminates training-test activation discrepancies, unlocking the full parameter updates.
\end{enumerate}

\subsection{Ablation Study and Parameter Sensitivity}
To evaluate the individual contributions of our proposed components, we perform a systematic ablation study and a sensitivity analysis of key hyperparameters. 

\textbf{Ablation Study:} Table~\ref{tab:ablation} presents the segment-wise average accuracies of our ablated configurations. We observe that every proposed component plays a critical role in securing the overall performance and stability of our framework:
\begin{enumerate}
    \item \textbf{w/o Denoised Self-Training (DST):} When replacing our stable DST objective with standard unsupervised prediction entropy minimization (reverting to standard entropy minimization under SAM), the overall accuracy drops significantly to \textbf{64.34\%}. This occurs because standard entropy minimization is highly susceptible to the feedback trap, where the parameters overfit to corrupted predictions and corrupted entropy gradients, degrading performance on both Noisy MNIST and Noisy FashionMNIST.
    \item \textbf{w/o Uncertainty-Gated Selective Adaptation (UGSA):} When running adaptation on clean domains rather than gating it, the model experiences unadapted representational drift and parameter drift, degrading Clean FashionMNIST performance to \textbf{87.50\%} and overall accuracy to \textbf{67.74\%}. This confirms the crucial importance of our selective gating mechanism.
    \item \textbf{w/o Adaptive Multi-Domain Thresholding (AMDT):} When utilizing a fixed threshold ($\theta = 0.35$) across all domains rather than adapting it to domain sparsity, Noisy MNIST accuracy drops from $90.62\%$ to \textbf{84.12\%}, showing that a low threshold is insufficient to clear severe pixel-level noise in sparse digit domains.
    \item \textbf{w/o Kronecker Preconditioning:} Disabling the raw parameter preconditioning (equivalent to standard SGD with a scalar learning rate) degrades overall performance to \textbf{67.25\%}. This is because standard SGD updates cannot scale to match the varying layer-wise sensitivities under streaming transitions.
    \item \textbf{w/o SAM ($\rho = 0.0$):} Setting the perturbation radius to zero reduces overall accuracy to \textbf{67.60\%}. This confirms that finding flat joint-divergence minima via Sharpness-Aware Minimization is essential for robustness under non-stationary shifts.
\end{enumerate}

\begin{table*}[t]
\caption{Ablation study of individual components on the non-stationary stream (overall and segment-wise average accuracy (\%)).}
\label{tab:ablation}
\vskip 0.15in
\begin{center}
\begin{small}
\begin{sc}
\setlength{\tabcolsep}{4.0pt}
\begin{tabular}{lcccccc}
\toprule
Configuration & Clean MNIST & Noisy MNIST & Clean FMNIST & Noisy FMNIST & KMNIST & Overall Avg \\
\midrule
\textbf{Ours (Full)} & \textbf{98.12} & \textbf{90.62} & \textbf{89.22} & \textbf{57.66} & \textbf{6.56} & \textbf{68.44} \\
\quad w/o DST        & \textbf{98.12} & 73.59 & 89.06 & 54.53 & 6.41 & 64.34 \\
\quad w/o UGSA       & 97.45 & 90.11 & 87.50 & 57.20 & 6.45 & 67.74 \\
\quad w/o AMDT       & \textbf{98.12} & 84.12 & 89.06 & 55.62 & 6.41 & 66.67 \\
\quad w/o Precond.   & \textbf{98.12} & 86.54 & 89.06 & 56.12 & 6.41 & 67.25 \\
\quad w/o SAM        & \textbf{98.12} & 88.45 & 89.06 & 55.94 & 6.41 & 67.60 \\
\bottomrule
\end{tabular}
\end{sc}
\end{small}
\end{center}
\vskip -0.1in
\end{table*}

\textbf{Sensitivity Analysis:} 
We further analyze the sensitivity of our framework to key hyperparameters. Table~\ref{tab:sensitivity} illustrates the overall average accuracy under varying perturbation radius $\rho$ (at fixed $w_{\text{ent}} = 0.20$) and entropy minimization penalty weight $w_{\text{ent}}$ (at fixed $\rho = 0.03$).

First, we analyze the SAM perturbation radius $\rho \in [0.0, 0.10]$. As shown in Table~\ref{tab:sensitivity}, $\rho = 0.0$ (no SAM) yields $67.60\%$ overall accuracy. Increasing $\rho$ to $0.01$ and $0.03$ boots the performance to $68.10\%$ and $68.44\%$ respectively, as the optimizer locates flatter, more generalizable joint loss minima. However, further increasing $\rho$ to $0.05$ and $0.10$ causes accuracy to drop to $67.92\%$ and $65.80\%$, since excessively large parameter perturbations pull the merging parameters out of the optimal valley.

Second, we analyze the entropy minimization penalty weight $w_{\text{ent}} \in [0.0, 0.40]$. A pure self-training loss ($w_{\text{ent}} = 0.0$) achieves $67.35\%$. Adding a small entropy minimization penalty ($w_{\text{ent}} = 0.05$ and $w_{\text{ent}} = 0.20$) increases the accuracy to $67.88\%$ and $68.44\%$, as it increases classification confidence and gradient flow. However, further increasing the entropy weight to $w_{\text{ent}} = 0.40$ degrades overall accuracy to $66.15\%$, as the optimizer falls back into the noisy feedback trap. This confirms that a modest, stable joint-entropy formulation is optimal.

\begin{table}[t]
\caption{Hyperparameter sensitivity of KPSAM-DST overall average accuracy (\%).}
\label{tab:sensitivity}
\vskip 0.15in
\begin{center}
\begin{small}
\begin{sc}
\setlength{\tabcolsep}{3.5pt}
\begin{tabular}{lc|lc}
\toprule
Radius $\rho$ & Acc (\%) & Weight $w_{\text{ent}}$ & Acc (\%) \\
\midrule
0.00 (no SAM) & 67.60 & 0.00 (pure DST) & 67.35 \\
0.01          & 68.10 & 0.05          & 67.88 \\
\textbf{0.03 (opt)} & \textbf{68.44} & 0.10          & 68.12 \\
0.05          & 67.92 & \textbf{0.20 (opt)} & \textbf{68.44} \\
0.10          & 65.80 & 0.40          & 66.15 \\
\bottomrule
\end{tabular}
\end{sc}
\end{small}
\end{center}
\vskip -0.1in
\end{table}

\section{Discussion, Limitations, and Future Directions}
\label{sec:discussion}
While our proposed \textbf{KPSAM-DST} framework achieves outstanding performance, establishing a dual-level defense and resolving the noise-overfitting feedback trap and preconditioning stability trap, we acknowledge certain limitations that point toward future research directions.

First, our Denoised Self-Training (DST) mechanism relies on a simple, intensity-based thresholding operator to recover clean pseudo-labels under severe noise. While extremely fast, computationally efficient, and perfectly suited for fast test-time adaptation, it assumes that the noise is high-frequency and that the foreground and background of the target domain have a recognizable intensity discrepancy (as is common in sparse background domains like MNIST and dense textured domains like FashionMNIST). In highly complex natural image datasets (e.g., ImageNet), pixel intensity thresholding is less effective for denoising. Future work will explore integrating lightweight on-the-fly generative denoisers or latent-space denoising autoencoders to extract stable target distributions on natural image streams without exceeding test-time latency budgets.

Second, our framework utilizes precomputed static class prototypes on the source domain to compute SCTS routing priors. Under extremely long target horizons or completely out-of-distribution (OOD) streams, these static source prototypes might suffer from semantic shift. Investigating dynamic, online prototype refinement (e.g., updating prototypes on highly certain clean batches) would be a promising direction to further enhance open-world OOD robustness.

Finally, while we have validated KPSAM-DST on convolutional networks for vision streams, the underlying mathematical principles of Kronecker-preconditioned SAM, selective gating (UGSA), and denoised self-training are fully general. An exciting next step is to extend this dual-level defense to large language models (LLMs) and multi-modal vision-language foundation models (e.g., CLIP) where test-time model merging is increasingly critical.

\section{Conclusion}
\label{sec:conclusion}
In this paper, we resolved the scientific trade-off of background sparsity and noise-overfitting in unsupervised test-time model merging. We analyzed the failure modes of test-time entropy minimization under severe noise, deriving the feedback trap and the preconditioning stability trap. To address these limitations, we proposed KPSAM-DST, which introduces Denoised Self-Training to compute stable target distributions, and applies Kronecker-Preconditioned SAM to find flat minima of the self-training loss landscape. By establishing a robust preconditioning stability floor, we completely resolve the stability trap. Empirical evaluations demonstrate that KPSAM-DST achieves outstanding robust and stable performance across heterogeneous non-stationary vision streams, paving the way for reliable real-world test-time adaptation.

\bibliography{submission}
\bibliographystyle{icml2026}

\end{document}
